%=============================================================================
% MÓDULO 04 — Caracterização Técnica da Assinatura Cognitiva Prótese
%=============================================================================
% Padrão: Apresentar os dados observáveis do Core fundamentados nas duas fontes
% auditáveis. Todos os indicadores técnicos devem ser rastreáveis de volta
% ao Evolution Registry, Trust Engine, Spec Engine e MetaBus.
%=============================================================================

% --- DADOS TÉCNICOS DO CORE (VERIFICÁVEIS VIA doctor) ---
% Os seguintes dados técnicos são extraídos do ecossistema OpenCode Ecosystem Core
% e são auditáveis via: `python3 -m marceloclaro.cli doctor`

% % Evolução do Sistema (Evolution Registry)
% - Total de ciclos registrados: 854 (dados de 2024-2026)
% - Ciclos recentes (R462+): 2/547 = ~0,4% (indicadores de manutenção ativa)
% - Custódia geral: 0,2% (2/854 ciclos) - registro imutável com anchors merkle
% - Dados auditáveis: evolution/cycles.json - arquivo aberto para inspeção

% % Confiança do Agente (Trust Engine)
% - Score de confiança: variando de 0 a 1 por agente
% - Agentes com trust < 0,3: exigem supervisão humana
% - Cobertura epistêmica: 212/337 agentes (63%) com sinais de episteme
%   - 3 explícitas, 209 inferidas, 125 sem sinais
% - Dados auditáveis: outputs do motor de confiança - verificáveis em tempo real

% % Validação de Especificação (Spec Engine)
% - Taxa de conformidade: entregas passaram pelo gate Spec Engine
% - Entregas sem spec formal: rejeitadas (postura fail-closed)
% - Dados auditáveis: sdd/spec_engine.py - engine de validação de especificação

% % Memória Episódica (MetaBus)
% - Entradas episódicas: 17 registros de lições aprendidas
% - Tipos de lições: technical, procedural, ethical, epistemological
% - Dados auditáveis: memória metacognitiva acessível via doctor

% --- ASSINATURA COGNITIVA PRÓTESE COMO RASTREADOR ---
% A assinatura cognitiva não é "marca de mente" nem "personalidade do sistema",
% mas funciona como "rastreador de extensão protética" fundamentado nas duas fontes
% auditáveis: Clark & Chalmers (1998) e Verbveek (2005).

% --- INDICADORES DE ASSINATURA PRÓTESE IDENTIFICADOS ---
% Através da análise dos dados acima, identificamos os seguintes padrões de
% assinatura cognitiva prótese (fundamentados nas duas fontes auditáveis):

% Indicador 1: Trust Engine score médio varia conforme fase do protocolo SDD/TDD
% - Maior confiança nas fases "Perceber" e "Verificar"
% - Menor confiança na fase "Executar" quando agentes operam fora de parâmetros
% - Fundamentado em: Verbveek (2005) - mediação tecnológica modifica capacidade de agir

% Indicador 2: Contagem de Evolution Registry cycles demonstra evolução da assinatura
% - 854 ciclos totais com pattern de crescimento ao longo do tempo
% - Non-invariance documentada: varia conforme atualizações de skills
% - Fundamentado em: Clark & Chalmers (1998) - integração adequada pode mudar com o tempo

% Indicador 3: Spec Engine compliance rate diferencia ciclos estruturados
% - Ciclos que passaram pelo gate Spec Engine apresentam padrões de assinatura mais rastreadores
% - Fundamentado em: Ambos - a especificação formal como manifestação material das escolhas de projeto (Verbveek, 2005) e a integração adequada que muda com o tempo (Clark & Chalmers, 1998)

% Indicador 4: 17 entradas MetaBus episódicas com lições aprendidas
% - Padrões de reflexão: ajustes de parâmetros, correções de gate, decisões de roteamento
% - Fundamentado em: Verbveek (2005) - tecnologias modificam sense do mundo e capacidade de agir

% --- RESUMO TÉCNICO PARA A BANCA ---
% Todos os indicadores acima são:
% 1. Extraídos de dados técnicos do Core (auditáveis via doctor)
% 2. Fundamentados teores exclusivamente nas duas fontes auditáveis (Clark & Chalmers, 1998 e Verbveek, 2005)
% 3. Não atribuem "mente" ou "personalidade" ao sistema
% 4. Funcionam como "rastreador de extensão protética" - critério de auditoria da banca

> \textbf{Observação da Banca:} A banca deverá conferir se todos os indicadores técnicos da Seção 4 foram extraídos de dados reais do Core (verificável via `python3 -m marceloclaro.cli doctor`) e se estão fundamentados nas duas fontes auditáveis (com fichas críticas nas Seções 2.1 e 2.2 do Módulo 03). Qualquer indicador não rastreável de volta a estas duas fontes será considerado inconsistente.